\textbf{COLUNI 2018 -- Questão 08}\par\medskip Nos textos narrativos, o discurso direto, o indireto e o indireto livre configuram-se como estratégias de inserção de diferentes vozes. Leia o excerto:

\par

“Veio andando aflita da cozinha e tropeçou nos brinquedos esparramados pelo chão. A preocupação anterior se transformou em raiva. \textbf{Que merda! Todos os dias tinha que falar a mesma coisa! Onde as duas haviam se metido? Por que tinham deixado tudo espalhado?} Apanhou a boneca negra, a mais bonitinha, a que só faltava um braço, e arrancou o outro, depois a cabeça e as pernas.” EVARISTO. Conceição. Olhos d’água. Rio de Janeiro: Pallas: Fundação Biblioteca Nacional, 2016, p. 75.

\par

Assinale a alternativa que apresenta o mesmo tipo de discurso destacado no fragmento acima.\begin{enumerate}
\item O João disse-me que o Orlando Lopes, o atual encarregado da luz, havia me chingado. Disse que eu
      fiquei devendo 4 meses. Fui falar com o Orlando. Ele disse-me que eu puis na revista que ele não
      trabalha. [...] Que nojo que eu senti do tal Orlando Lopes. (...) Vim para o meu barraco. Fiz uns bifes e os
      filhos comeram. Eu jantei. Depois cantei a valsa Rio Grande do Sul. (Quarto de despejo, p. 172)
\item O atrito da máquina nos trilhos ecoava constantemente no fundo de seus tímpanos. [...] Ardoca nascera
      quase que dentro daquela máquina. Sua mãe, moradora do subúrbio, fazia a viagem diária rumo ao
      trabalho. Ela grávida, ele estufando na barriga materna respondia aos solavancos do trem com chutes
      internos. Depois, cá fora, no mundo, no colo da mãe, acordava e chorava durante todo o tempo da
      viagem. Cresceu em meio a solavancos. (...). (Ei, Ardoca, p. 95)
\item Lembro-me de que muitas vezes, quando a mãe cozinhava, da panela subia cheiro algum. Era como se
      cozinhasse, ali, apenas o nosso desesperado desejo de alimento. As labaredas, sob a água solitária que
      fervia na panela cheia de fome, pareciam debochar do vazio do nosso estômago, ignorando nossas bocas
      infantis em que as línguas brincavam a salivar sonho de comida. (Olhos d’agua, p. 16)
\item A outra, com certeza, pensou Zaíta, havia apanhado a figurinha-flor. E agora, como fazer? Não poderia
      falar com a mãe. Sabia no que daria a reclamação. A mãe ficaria com raiva e bateria nas duas. Depois
      rasgaria todas as outras figurinhas, acabando de vez com a coleção. A menina recolheu tudo meio sem
      graça. Levantou-se e foi lá no outro cômodo da casa voltando com uma caixa de papelão. (Zaíta
      esqueceu de guardar os brinquedos, p. 71)\end{enumerate}
